\documentclass[12pt]{article}

% Настройка русского языка и шрифтов
\usepackage[utf8]{inputenc}
\usepackage[russian]{babel}
\usepackage[T2A]{fontenc}

% Математические пакеты
\usepackage{amsmath, amssymb, amsfonts}

% Настройка страницы
\usepackage{geometry}
\geometry{a4paper, margin=0.8in}

% Другие полезные пакеты
\usepackage{hyperref}
\usepackage{booktabs}
\usepackage{enumitem}
\usepackage{longtable}
\usepackage{graphicx}
\usepackage{xcolor}

% Определение стилей для комментариев
\definecolor{commentgray}{rgb}{0.4,0.4,0.4}
\definecolor{antigravityblue}{rgb}{0.0, 0.45, 0.74}

\newcommand{\commentary}[1]{{\small\itshape\color{commentgray} #1}}
\newcommand{\antigravity}[1]{{\small\color{antigravityblue}\textbf{Antigravity:} #1}}

% Настройка ссылок
\hypersetup{
    colorlinks=true,
    linkcolor=blue,
    filecolor=magenta,      
    urlcolor=cyan,
    pdftitle={Глубокий разбор: Метрические методы},
}

\title{Глубокий разбор: Метрические методы обучения}
\author{Твой персональный гид по ML}
\date{\today}

\begin{document}

\subsection{Алгоритм k-ближайших соседей (k-NN)}

\begin{itemize}
    \item \textbf{Математическая формулировка}: Пусть задана обучающая выборка $X = \{(x_i, y_i)\}$. Для нового объекта $x$ мы находим $k$ ближайших к нему объектов из $X$: $x_{(1)}, \dots, x_{(k)}$.
    \item \textbf{Классификация}: Ответ — наиболее часто встречающийся класс среди соседей (мода): $a(x) = \text{argmax}_{y \in Y} \sum_{j=1}^k [y_{(j)} = y]$.
    \item \textbf{Регрессия}: Ответ — среднее арифметическое ответов соседей: $a(x) = \frac{1}{k} \sum_{j=1}^k y_{(j)}$.
    \item \textbf{Параметр $k$}: 
    \begin{itemize}
        \item При \textbf{малых $k$} модель имеет \textbf{низкое смещение (bias)}, но \textbf{высокий разброс (variance)}.
        \item При \textbf{больших $k$} модель имеет \textbf{высокое смещение}, но \textbf{низкий разброс}.
    \end{itemize}
\end{itemize}

\commentary{Выбор $k$ — это классический поиск компромисса между недообучением и переобучением. Обычно $k$ подбирают на кросс-валидации.}

\subsection{Метрики и функции расстояния}

\begin{itemize}
    \item \textbf{Метрики Минковского ($L_p$)}: $\rho(x, z) = \left( \sum_{i=1}^d |x_i - z_i|^p \right)^{1/p}$. 
    \begin{itemize}
        \item $L_1$ (Манхэттенское): сумма модулей разностей. Менее чувствительна к выбросам.
        \item $L_2$ (Евклидово): стандартное расстояние по прямой.
    \end{itemize}
    \item \textbf{Косинусное сходство}: $\cos(\theta) = \frac{\langle x, z \rangle}{\|x\| \|z\|}$. Расстояние определяется как $1 - \cos(\theta)$. Измеряет угол между векторами, игнорируя их длину.
    \item \textbf{Расстояние Жаккара}: $1 - \frac{|A \cap B|}{|A \cup B|}$. Используется для сравнения множеств (стандарт для бинарных признаков).
\end{itemize}

\commentary{Помните, что перед использованием $L_p$ метрик данные \textbf{обязательно} нужно масштабировать, иначе признак с большим диапазоном значений «поглотит» все остальные.}

\subsection{Модификации: Взвешенный k-NN}

\begin{itemize}
    \item \textbf{Веса по рангу}: Вес $j$-го соседа зависит от его номера: $w_j = \frac{k+1-j}{k}$.
    \item \textbf{Веса от расстояния}: $w_i = f(\rho(x, x_i))$. Позволяет доверять более близким объектам сильнее.
    \item \textbf{Ядерное сглаживание}: Вес задается через ядро $K$: $w_i = K\left(\frac{\rho(x, x_i)}{h}\right)$. 
    \item \textbf{Примеры ядерных функций} (в порядке увеличения их гладкости): 
    \begin{itemize}
        \item Прямоугольное: $K(z) = \frac{1}{2} [|z| \le 1]$.
        \item Треугольное: $K(z) = (1 - |z|) [|z| \le 1]$.
        \item Епанечникова: $K(z) = \frac{3}{4} (1 - z^2) [|z| \le 1]$.
        \item Гауссовское: $K(z) = \frac{1}{\sqrt{2\pi}} e^{-\frac{1}{2}z^2}$ (бесконечно гладкое).
    \end{itemize}
\end{itemize}

\subsection{Ядерная регрессия (Формула Надарая-Ватсона)}

\begin{itemize}
    \item \textbf{Принцип}: Это обобщение k-NN, где предсказание — это взвешенное среднее ответов \textbf{всех} объектов обучающей выборки: 
    $a(x) = \frac{\sum_{i=1}^n y_i K\left(\frac{\rho(x, x_i)}{h}\right)}{\sum_{i=1}^n K\left(\frac{\rho(x, x_i)}{h}\right)}$.
    \item \textbf{Ширина окна $h$}: 
    \begin{itemize}
        \item Малое $h$: веса быстро затухают (риск переобучения).
        \item Большое $h$: прогноз становится слишком гладким (риск недообучения).
    \end{itemize}
\end{itemize}

\section{Оптимизация поиска ближайших соседей}

\begin{itemize}
    \item \textbf{k-d деревья}: Бинарное дерево, разбивающее пространство вдоль осей. Эффективно при $d < 20$.
    \item \textbf{Приближенные методы (ANN)}: Согласие найти "достаточно близкого" соседа для ускорения в тысячи раз.
    \item \textbf{Locality-sensitive hashing (LSH)}: Хеш-функция переводит близкие объекты в один бакет. Вероятность коллизии для близких объектов должна быть высокой, а для далёких — низкой.
    \item \textbf{Hierarchical Navigable Small World (HNSW)}: На данных строится многоуровневый NSW-граф, где между любыми двумя точками существует короткий путь. Исходный граф — нулевой слой, а каждый следующий содержит всё меньше вершин, позволяя быстро "перепрыгивать" по пространству.
\end{itemize}

\antigravity{В реальном продакшене (Faiss, HNSW) почти всегда используются именно ANN, так как точный поиск на миллионах объектов невозможен.}

\end{document}
